\section{API dead ends}

\subsection{Case Study 1: Metaculus 403 API Deprecation and Tier Lock (v0.5.1)}

\subsubsection*{Контекст}
Metaculus служит источником структурированных прогнозов сообщества в stage 1 (ForesightCollector). Их API интегрирована в \texttt{src/data\_sources/foresight.py} с 2026-03-28. На production сервере v0.5.1 все запросы к Metaculus возвращали HTTP 403 Forbidden, но локально тесты проходили. Это свидетельствовало о скрытом breaking change в API.

\subsubsection*{Что пробовали}
В диагностическом скрипте (\texttt{tasks/research/metaculus\_api\_fix.md}) проверили несколько гипотез:
\begin{enumerate}
\item GET requests требуют авторизацию (было предположение: read-доступ свободный)
\item Параметры изменили названия между версиями (`status` vs `statuses`)
\item Endpoint \texttt{/api2/questions/} всё ещё работает как документировано
\end{enumerate}

\subsubsection*{Почему не сработало}
Две причины одновременно:

\textbf{(1) API v2 fully deprecated.} Примерно в Q3--Q4 2024 Metaculus перенёсли все endpoints с \texttt{/api2/*} на \texttt{/api/posts/*}. Документация не синхронизирована. Endpoint \texttt{/api2/questions/} остался работающим 10--15 месяцев, создав впечатление обратной совместимости. На самом деле, поддержка заканчивается.

\textbf{(2) Все endpoints требуют авторизацию.} Историческая норма: GET-запросы к публичным API бесплатные без токена. Metaculus изменили политику --- даже чтение вопросов требует `Authorization: Token <token>` заголовок. Возможно, это было сделано для rate-limiting по пользователю.

\subsubsection*{Что сделали}
Fix реализован в v0.5.1 (коммит в CHANGELOG.md, стр. 15):

\begin{enumerate}
\item Конструктор MetaculusClient принимает опциональный параметр \texttt{token: str = ""}. Если указан, заголовок \texttt{Authorization: Token {token}} добавляется в каждый запрос.
\item Endpoint изменён с \texttt{/api2/questions/} на \texttt{/api/posts/}.
\item Параметры переименованы:
    \begin{itemize}
    \item `status` $\to$ `statuses`
    \item \texttt{resolve\_time\_\_gt} $\to$ \texttt{scheduled\_resolve\_time\_\_gt}
    \item \texttt{resolve\_time\_\_lt} $\to$ \texttt{scheduled\_resolve\_time\_\_lt}
    \end{itemize}
\item Парсинг ответа переписан. Новый формат:
\begin{verbatim}
post['question']['aggregations']['recency_weighted']['latest']['centers'][0]
\end{verbatim}
вместо старого:
\begin{verbatim}
post['community_prediction']['full']['q2']
\end{verbatim}
\item Токен генерируется на https://www.metaculus.com/aib/ (бесплатный, не истекает).
\end{enumerate}

v0.9.4 (CHANGELOG, стр. 83) Metaculus отключён в production (\texttt{\_fetch\_metaculus()} возвращает \texttt{[]} если client=None). Причина: access tier ограничен --- нам нужен BENCHMARKING tier для надёжного доступа.

\subsubsection*{Урок}
\textbf{Никогда не предполагай обратную совместимость третьего-сторонней API.} Даже если старый endpoint работает 10+ месяцев, это может быть legacy layer, а не guarantee.
\begin{enumerate}
\item Подписывайся на изменения API через официальные каналы (RSS, email newsletter, GitHub releases).
\item При интеграции внешней API, сохраняй последнюю документацию локально (\texttt{tasks/research/metaculus\_api\_fix.md} служит reference).
\item Зафиксируй rate limit --- Metaculus: 120 req/min (безопасно), без опубликованного лимита. Всегда консервативен.
\item При auth-требовании, явно документируй, как получить credentials, и сделай их конфигурируемыми из окружения.
\end{enumerate}

---

\subsection{Case Study 2: GDELT Cyrillic Query Crash (pre-v0.5.1, v0.9.4)}

\subsubsection*{Контекст}
ForesightCollector вызывает GDELT DOC 2.0 API с query вида \emph{``news forecast 2026-03-29 ТАСС''} для сбора статей по русским СМИ. На production сервере с v0.9.2 и v0.9.3 это вызывало JSON parse error: ElasticSearch backend за GDELT возвращал HTML вместо JSON.

\subsubsection*{Что пробовали}
В \texttt{tasks/research/gdelt\_api\_fix.md} и \texttt{tasks/research/gdelt\_api.md} проведён анализ:
\begin{enumerate}
\item GDELT query-язык поддерживает Unicode кириллицу?
\item API имеет опции для фильтрации по языку источника?
\item Как парсить 403 HTML-ответ, чтобы не падать?
\item Есть ли более структурированный источник событий (не pure search API)?
\end{enumerate}

\subsubsection*{Почему не сработало}
Три причины одновременно:

\textbf{(1) Кириллица в query.} GDELT индексирует только английские переводы и метаданные. Когда ForesightCollector отправляет \texttt{query="news forecast 2026-03-29 ТАСС"}, ElasticSearch не может обработать кириллицу в текущей конфигурации. Вместо JSON error, API возвращает HTML-страницу (5xx error page).

\textbf{(2) Будущие даты бессмысленны.} GDELT основан на rolling window ~ 3 месяца в прошлое. Query с датой 2026-03-29 (из futuer) не вернёт никаких статей. API не валидирует дату --- просто возвращает пустой результат (или ошибку).

\textbf{(3) Отсутствие защиты от null в `articles`.} Когда API испытывает нагрузку, может вернуть \texttt{{"articles": null}} вместо \texttt{{"articles": []}}. Код \texttt{for article in data.get("articles", [])} упадёт с TypeError, потому что `None` не iterable.

\subsubsection*{Что сделали}
Fix реализован в два этапа: Content-Type guard + null-safe articles в pre-v0.5.1 (commit 987260c), Cyrillic stripping в v0.9.4 (commit 3658531):

\begin{enumerate}
\item \textbf{Content-Type проверка.} Перед парсингом JSON:
\begin{verbatim}
if "text/html" in response.headers.get("content-type", ""):
    logger.warning("GDELT returned HTML: %s", response.text[:200])
    return []
\end{verbatim}

\item \textbf{Null-safe articles.} Замена:
\begin{verbatim}
for article in (data.get("articles") or []):
\end{verbatim}
вместо \texttt{data.get("articles", [])}, потому что latter не защищает от `null`.

\item \textbf{Английский query.} ForesightCollector теперь генерирует query на английском: `"news forecast"` вместо `"Новостной прогноз"`. Language-specific фильтрация делается через GDELT operators (ниже).

\item \textbf{Language и country operators.} Документация в \texttt{tasks/research/gdelt\_api.md} (стр. 8--9) показывает:
\begin{verbatim}
sourcelang:russian + sourcecountry:RS
\end{verbatim}
(RS = FIPS код России, не ISO RU). Это фильтрует статьи по языку источника и стране издателя без кириллицы в query.
\end{enumerate}

\subsubsection*{Урок}
\textbf{Никогда не доверяй charset-поддержке без явного тестирования.} Даже если API документация не упоминает ограничения, всегда проверяй граничные случаи (non-ASCII characters, future dates, null values).

\begin{enumerate}
\item Запрашиваемые языки обрабатывай через фильтр (e.g. `sourcelang:` operator), а не через query text.
\item Проверяй `content-type` заголовок перед парсингом JSON --- это спасает от HTML error pages.
\item Защищайся от null values: \texttt{(data.get("field") or [])} вместо \texttt{data.get("field", [])}.
\item Rate limit на GDELT: ~1 req/sec, без опубликованного лимита. Добавь exponential backoff на 429.
\end{enumerate}

---

\subsection{Case Study 3: Polymarket camelCase Parameter Mismatch (pre-v0.5.1)}

\subsubsection*{Контекст}
Polymarket Gamma API используется в ForesightCollector для fetch markets. Код в \texttt{src/data\_sources/foresight.py} (строка 165) отправлял:
\begin{verbatim}
GET /markets?order=volume_24hr
\end{verbatim}

На production это возвращало HTTP 422 Unprocessable Entity. Локально в dry-run не воспроизводилось (возможно, из-за кэша или другого контекста).

\subsubsection*{Что пробовали}
Live-запросы (2026-03-28) с разными вариантами параметра:
\begin{enumerate}
\item \texttt{order=volume\_24hr} (snake\_case, как документировано в algunos examples) $\to$ 422
\item \texttt{order=volume24hr} (camelCase, как в других частях Polymarket API) $\to$ 200
\item \texttt{order=volume} (generic, не специфично) $\to$ 200
\end{enumerate}

Вывод: параметр требует camelCase, но никаких доков не было.

\subsubsection*{Почему не сработало}
\textbf{Inconsistency между endpoints.} Polymarket имеет несколько API endpoints:
\begin{itemize}
\item \texttt{/markets} --- принимает camelCase (\texttt{volume24hr}, \texttt{liquidityNum}, \texttt{endDate})
\item \texttt{/events} --- принимает snake\_case (\texttt{volume\_24h}, \texttt{liquidity\_num})
\end{itemize}

Это расхождение нигде не задокументировано. Разработчик (справедливо) предположил, что snake\_case универсален. На самом деле, \texttt{/markets} настаивает на camelCase.

\subsubsection*{Что сделали}
Fix реализован в pre-v0.5.1 (commit 987260c, 2026-03-28):

Изменение в \texttt{src/data\_sources/foresight.py}:
\begin{verbatim}
# БЫЛО:
"order": "volume_24hr"

# СТАЛО:
"order": "volume24hr"
\end{verbatim}

Документация в \texttt{tasks/research/polymarket\_gamma\_api\_fix.md} (стр. 26--46) содержит полный reference валидных параметров для \texttt{/markets} endpoint.

\subsubsection*{Урок}
\textbf{API inconsistency --- это bug в дизайне API, не вашего кода.} Однако, когда ты интегрируешь, нужно быть готов к такому.

\begin{enumerate}
\item Всегда тестируй все endpoints той же API live перед production deploy.
\item Не полагайся на consistency между похожими endpoints --- каждый может иметь свою convention.
\item Создай локальный reference документ для каждого внешнего API (как \texttt{tasks/research/}) с примерами.
\item При 422 Unprocessable Entity, форматируй запрос как JSON и проверь каждый параметр через документацию.
\end{enumerate}

---

\subsection{Case Study 4: Reuters RSS Feed Deprecation (v0.9.4)}

\subsubsection*{Контекст}
Исторически GLOBAL\_RSS\_FEEDS включал Reuters фиды:
\begin{verbatim}
https://feeds.reuters.com/reuters/...
\end{verbatim}

На production после 2026-03-28 эти URLs возвращали 404, не обновляясь. Pipeline падал на stage 1 при попытке fetch news signals через эти dead feeds.

\subsubsection*{Что пробовали}
Проверили несколько гипотез:
\begin{enumerate}
\item Reuters сменили хост (e.g. reuters.com/news vs feeds.reuters.com)?
\item Feed требует авторизацию?
\item Хост DNS разрешается, но HTTP 404 возвращается?
\end{enumerate}

Итог: \texttt{feeds.reuters.com} закрыт полностью (HTTPS SSL certificate expired, не поддерживается). Сервис закрыт с 2020 года.

\subsubsection*{Почему не сработало}
Reuters мигрировала на внутренний ReutersConnect API (требует авторизация, недоступна для public). Public RSS feeds больше не поддерживаются.

\subsubsection*{Что сделали}
Удалены все Reuters feeds из \texttt{src/agents/collectors/news\_scout.py} (GLOBAL\_RSS\_FEEDS). Альтернативные источники (BBC, AP, Bloomberg) остаются.

\subsubsection*{Урок}
\textbf{RSS feeds --- живые ресурсы. Периодически проверяй.} Крупные издатели могут закрыть public feeds в любой момент.

\begin{enumerate}
\item Добавь health-check для RSS feeds (попытайся fetch за последнюю неделю).
\item Имей fallback источники --- Reuters закрыт, но BBC, AP, Bloomberg, свежие.
\item Не полагайся на одного издателя для news signals.
\item Отслеживай 404/410 в логах --- это сигнал, что feed мёртв.
\end{enumerate}

---

\section{Архитектурные отказы}

\subsection{Case Study 5: YandexGPT Stub (v0.8.0)}

\subsubsection*{Контекст}
v0.7.0 содержал условную логику для выбора LLM провайдера: если \texttt{OPENROUTER\_API\_KEY} отсутствует, fallback на YandexGPT (\texttt{ya\_gpt\_pro} или \texttt{ya\_gpt\_lite}). YandexGPT требует российский IP (или VPN) и YandexCloud аккаунт.

\subsubsection*{Почему не сработало}
\begin{enumerate}
\item Production сервер в Yandex Cloud (213.165.220.144) имеет доступ к YandexGPT, но интеграция бросала \texttt{NotImplementedError} при каждом вызове (никогда не была реализована дальше stub'а).
\item Приоритизация провайдера была жёсткая: если YandexGPT задан, OpenRouter не используется. Это создавало single point of failure.
\item Документация prompts отсутствовала специфична для YandexGPT (модель имеет своё поведение JSON parsing).
\item Обслуживание двух параллельных стеков (OpenRouter + YandexGPT) удваивает complexity.
\end{enumerate}

\subsubsection*{Что сделали}
v0.8.0 (CHANGELOG, стр. 231): YandexGPT полностью удален. Три задачи, которые раньше использовали YandexGPT (\texttt{style\_generation}, \texttt{style\_generation\_ru}, \texttt{quality\_style}), переведены на OpenRouter с \texttt{anthropic/claude-sonnet-4}.

\subsubsection*{Урок}
\textbf{Не строй архитектуру на наличии двух провайдеров, если один из них не работает.} Single LLM provider (OpenRouter) лучше, чем fallback на неработающего second.

\begin{enumerate}
\item Выбери один primary провайдер и придерживайся.
\item Если нужна redundancy, выбери два провайдера, оба fully tested и integrated.
\item Не добавляй условную логику ``если API key отсутствует, используй другого провайдера'' --- это скрывает problems.
\end{enumerate}

---

\subsection{Case Study 6: Non-Existent Preset Sonnet 4.6 (v0.9.4)}

\subsubsection*{Контекст}
v0.9.3 содержал 3 пресета в \texttt{src/config.py}:
\begin{enumerate}
\item Light (Gemini Flash)
\item Standard (Claude Sonnet 4.6)
\item Opus (Claude Opus 4.6)
\end{enumerate}

Standard пресет был фасадом для модели \texttt{claude-sonnet-4.6}, которой нет в OpenRouter pricing. Результат: при выборе Standard, pipeline падал на первом LLM call с ошибкой ``model not found''.

\subsubsection*{Почему не сработало}
\begin{enumerate}
\item Confusion между версиями Sonnet. OpenRouter поддерживает \texttt{claude-3.5-sonnet}, но не \texttt{claude-sonnet-4.6} (и не существует в природе).
\item Пресеты не были протестированы полностью (dry-run использовал флаги, не UI-выбор).
\item Документация была неактуальна (старые имена моделей).
\end{enumerate}

\subsubsection*{Что сделали}
v0.9.4 (CHANGELOG, стр. 85): Standard пресет удален. Оставлены 2 пресета: Light (Gemini 2.5 Flash) и Opus (Claude Opus 4.6). UI обновлена.

\subsubsection*{Урок}
\textbf{Валидируй model names до деплоя.} Если пресет ссылается на модель, проверь, что она существует в pricing.

\begin{enumerate}
\item Добавь CI/CD step, который тестирует каждый пресет с real API call.
\item Держи актуальный список доступных моделей на каждом провайдере.
\item Лучше иметь 1 пресет, который работает, чем 3, из которых 2 сломаны.
\end{enumerate}

---

\subsection{Case Study 7: Dark Mode Complexity (v0.8.0)}

\subsubsection*{Контекст}
v0.7.0 и ранее содержали CSS для dark mode (Pico.css имеет встроенную поддержку). Было HTML-переключение `<button>` для toggle. Это добавило:
\begin{enumerate}
\item CSS сложность (duplication для light/dark палеток)
\item JS сложность (localStorage для preference, system preference detection)
\item Testing complexity (тестировать оба режима)
\item User confusion (не все пользователи знают, что toggle существует)
\end{enumerate}

\subsubsection*{Почему не сработало}
\begin{enumerate}
\item Нет пользовательского спроса на dark mode (feedback не указывал на это).
\item Дизайн был субоптимален --- цвета светлого режима были тщательно выбраны (OKLCH, контрастность, доступность). Dark mode версия была скопирована без адаптации.
\item Maintenance burden: каждое изменение цвета требовало обновления двух палеток.
\end{enumerate}

\subsubsection*{Что сделали}
v0.8.0 (CHANGELOG, стр. 231): Dark mode удален. Оставлена только светлая тема (OKLCH палетка, оптимизированная для контрастности и доступности).

\subsubsection*{Урок}
\textbf{Не добавляй features без спроса. YAGNI.} Даже если feature ``nice to have'', это не стоит complexity.

\begin{enumerate}
\item Собирай feedback перед добавлением UI features.
\item Если feature не используется, удали её.
\item Лучше одна хорошо-сделанная тема, чем две плохо-сделанные.
\end{enumerate}

---

\subsection{Case Study 8: Миграция Pico.css на Tailwind CSS (v0.8.0)}

\subsubsection*{Контекст}
Начальная версия UI использовала Pico.css --- classless CSS-фреймворк (минимальная стилизация без классов). Для прототипов этого было достаточно: формы, таблицы и типографика выглядели приемлемо из коробки.

\subsubsection*{Почему не сработало}
С развитием продукта требования к UI выросли за пределы возможностей classless-фреймворка:
\begin{enumerate}
\item \textbf{OKLCH-палитра.} Pico.css работает на HSL/hex, не поддерживает OKLCH цветовое пространство, необходимое для perceptually uniform палитры дизайн-системы.
\item \textbf{Кастомные компоненты.} Dashboard с прогнозами, карточки результатов, timeline --- ни один из этих элементов не покрывается classless CSS.
\item \textbf{Responsive layout.} Pico.css предоставляет базовый responsive через media queries, но нет утилит для grid/flex раскладки dashboard'ов.
\item \textbf{Анимации.} Переходы между состояниями (loading, success, error) требуют кастомного CSS, что полностью нивелирует преимущество ``zero-class'' подхода.
\end{enumerate}

\subsubsection*{Что сделали}
v0.8.0 (CHANGELOG, стр. 231): полная миграция на Tailwind CSS v4 с PostCSS build pipeline (commit ae12705):

\begin{enumerate}
\item Tailwind CSS v4.2.2 с \texttt{@theme} config заменил Pico.css.
\item Реализована дизайн-система Impeccable с 17 JS-referenced \texttt{fn-*} компонентами.
\item Типографика: Newsreader (заголовки) / Source Sans 3 (body) / JetBrains Mono (код) через Google Fonts.
\item PostCSS build: \texttt{input.css} $\to$ \texttt{tailwind.css} с watch mode для разработки.
\end{enumerate}

\subsubsection*{Урок}
\textbf{Classless CSS-фреймворки --- ловушка для всего, что сложнее landing page.} Если UI потребует кастомных компонентов, утилитарный фреймворк нужно выбирать сразу.

\begin{enumerate}
\item Оценивай потолок UI-требований на старте проекта, а не по текущим потребностям.
\item Стоимость миграции CSS растёт нелинейно с количеством шаблонов.
\item Утилитарный фреймворк (Tailwind) масштабируется с проектом; classless --- нет.
\end{enumerate}

---

\section{Критические баги}

\subsection{Case Study 9: Temporal Leak in Walk-Forward Evaluation (v0.9.2)}

\subsubsection*{Контекст}
Phase 4 (v0.9.2) реализовала walk-forward evaluation для валидации informed consensus. Модель использовала pre-aggregated позиции бетторов (из старого JSON-датасета) за каждый рынок. При итерации по 22 фолдам walk-forward каждый фолд имеет temporal cutoff $T$ (``используй только данные до даты $T$'').

Проблема: pre-aggregated позиции содержали трейды со средней датой $T + 30$ дней. Это means информационная утечка из будущего --- модель видит сигналы, которые не могла видеть реальная система в момент времени $T$.

\subsubsection*{Что пробовали}
\begin{enumerate}
\item Проверить, что cutoff $T$ корректен? Да, cutoff применялась правильно.
\item Проверить, что фолды не перекрываются? Да, non-overlapping 60-day windows.
\item Где источник утечки? Traced в pre-aggregated \texttt{avg\_position} для каждого bettor-market пары.
\end{enumerate}

\subsubsection*{Почему не сработало}
\textbf{Design issue.} Pre-aggregated positions были computed один раз на весь датасет:
\begin{verbatim}
avg_position = SUM(price $\times$ quantity) / SUM(quantity) for all trades
\end{verbatim}

Это глобальное среднее, без временного измерения. При walk-forward, код использовал то же самое среднее для всех фолдов, забывая, что медиана датирована будущим.

\subsubsection*{Что сделали}
v0.9.2 (CHANGELOG, стр. 119--122): Новый скрипт \texttt{scripts/duckdb\_build\_bucketed.py} переагрегирует 470M raw trades в bucketed parquet:

\begin{enumerate}
\item Временные бакеты: 30-дневные окна (e.g., 2024-01-01 до 2024-01-31).
\item Для каждого (bettor, market, time\_bucket): сохрани \texttt{weighted\_price\_sum} и \texttt{total\_usd} (частичные суммы).
\item Walk-forward теперь может вычислить: $\text{avg\_position\_as\_of\_T} = \frac{\sum \text{weighted\_price\_sum}}{\sum \text{total\_usd}}$ где $\text{time\_bucket} \le T$.
\item Zero data loss: 470M trades $\to$ 2.4GB bucketed parquet с predicate pushdown.
\end{enumerate}

Result (CHANGELOG, стр. 127): 22/22 фолда $\text{BSS} > 0$ с robust mean $\text{BSS} = +0.127$.

Анализ утечки (CHANGELOG, стр. 130): leaked $\text{BSS} = +0.092$ vs clean $\text{BSS} = +0.117$ на тех же фолдах. Leak добавлял шум, не сигнал.

\subsubsection*{Урок}
\textbf{Temporal cutoff требует временного измерения в данных, не глобальное агрегирование.} Это критично для любого walk-forward evaluation.

\begin{enumerate}
\item Никогда не используй pre-aggregated data для temporal validation.
\item Bucketed/timestamped aggregates позволяют point-in-time queries: ``какое было среднее на 2024-01-15?''
\item Всегда проверяй: есть ли trades/signals с датой > cutoff T? Если да, это leak.
\item Unit test: add explicit check that \texttt{as\_of\_date} < cutoff date for all test data.
\end{enumerate}

---

\subsection{Case Study 10: conditionId Mismatch (v0.9.3)}

\subsubsection*{Контекст}
Polymarket имеет дуальную структуру: каждый рынок имеет two IDs:
\begin{itemize}
\item `id` (numeric, e.g., 123456) --- локальный в Gamma API marketplace
\item `conditionId` (CTF hex hash, e.g., 0x7e2a...) --- глобальный в CLOB и Data API
\end{itemize}

Phase 6 (v0.9.5, но bug был в v0.9.3) интегрировала Data API для fetch real trades. ForesightCollector пытается обогатить Gamma markets с informèd consensus signals из trade data.

Проблема: \texttt{src/data\_sources/foresight.py} использовал \texttt{market.id} (numeric) как join key, а \texttt{src/inverse/loader.py} (который загружал исторические trades) использовал `conditionId` (hash). 99% сигналов молчал отпадал, потому что не было match'а.

\subsubsection*{Что пробовали}
\begin{enumerate}
\item Посмотреть, что возвращает Polymarket API? Оба поля возвращаются.
\item Какое поле используется в других частях кода? Нашли inconsistency.
\item Добавить логирование для match rate? Обнаружили 0% matches.
\end{enumerate}

\subsubsection*{Почему не сработало}
\textbf{Нет documentation in code comments.} Ни один из двух IDs не был задокументирован с пояснением, что один локальный, другой глобальный. Разработчик справедливо выбрал первый встреченный `id` поле.

\subsubsection*{Что сделали}
v0.9.3 (CHANGELOG, стр. 96--97):

\begin{enumerate}
\item \texttt{src/data\_sources/foresight.py}: extract `conditionId` вместо `id`.
\item \texttt{src/agents/collectors/foresight\_collector.py}: use `conditionId` как join key для matching с trade data.
\item Добавлены comments:
\begin{verbatim}
# conditionId is globally unique (CLOB + Data API)
# id is local to Gamma API (marketplace-specific)
# Always join on conditionId, not id.
\end{verbatim}
\end{enumerate}

\subsubsection*{Урок}
\textbf{Когда third-party API имеет несколько IDs, документируй purpose каждого.}

\begin{enumerate}
\item Add comment в код: ``// market.id: local (Gamma), market.conditionId: global (CLOB)''.
\item Test: fetch market через оба endpoints, verify оба ID'а return.
\item Unit test: join на wrong ID должен вернуть 0 matches (catch этот bug).
\item Добавь assertion: \texttt{assert len(matches) > 0, "No markets matched! Check join keys."}
\end{enumerate}

---

\subsection{Case Study 11: Date Serialization Crash (v0.9.5)}

\subsubsection*{Контекст}
v0.9.5 обновила timeline schemas (Phase 6). \texttt{predicted\_timeline} содержит \texttt{predicted\_date} и \texttt{target\_date} поля. Pipeline прошла all 9 stages успешно, но при попытке сохранить результат в SQLite, worker упал с:

\begin{verbatim}
TypeError: Object of type date is not JSON serializable
\end{verbatim}

Результат не был сохранён --- user смотрит на пустую страницу despite 40 минут computation (\$5--15 в cost).

\subsubsection*{Что пробовали}
\begin{enumerate}
\item Проверить, что dates парсятся корректно? Да, они \texttt{datetime.date} объекты.
\item Где serialization происходит? В worker при \texttt{model\_dump()} перед `INSERT`.
\item Почему \texttt{model\_dump()} не обрабатывает date объекты? По дефолту, \texttt{model\_dump()} не имеет \texttt{mode="json"}, который преобразует non-JSON types.
\end{enumerate}

\subsubsection*{Почему не сработало}
\textbf{Pydantic v2 breaking change.} В Pydantic v2, \texttt{model\_dump()} требует явный \texttt{mode="json"} для JSON serialization. По дефолту, это просто возвращает Python объекты. Тогда \texttt{json.dumps()} fails на `date` объектах.

\subsubsection*{Что сделали}
v0.9.5 (CHANGELOG, стр. 35--36):

\begin{enumerate}
\item Добавлены \texttt{@field\_serializer} decorators на schemas с date полями:
\begin{verbatim}
@field_serializer('predicted_date', 'target_date', 'generated_at')
def serialize_dates(self, value):
    if isinstance(value, date):
        return value.isoformat()
    return value
\end{verbatim}

\item Альтернативно: используй \texttt{model\_dump(mode="json")} везде, где нужна JSON serialization.
\item Добавлены unit tests для timeline serialization.
\end{enumerate}

\subsubsection*{Урок}
\textbf{Serialization должна быть unit-tested. Никогда не предполагай, что \texttt{model\_dump()} даст тебе JSON-ready dict.}

\begin{enumerate}
\item Add test: \texttt{model\_dump() $\to$ json.dumps() $\to$ json.loads()} round-trip.
\item При добавлении новых fields типа date/datetime, immediately add serializer.
\item Используй \texttt{model\_dump(mode="json")} в production code (более явно).
\item Помни: Pydantic v2 требует explicit `mode` для JSON types.
\end{enumerate}

---

\subsection{Case Study 12: PromptParseError Silently Dropped (v0.5.2--v0.9.4)}

\subsubsection*{Контекст}
v0.7.0 agent (EventTrendAnalyzer) отправляет JSON request на LLM. LLM (Gemini Flash) может вернуть truncated JSON (если finish\_reason=``length''). Код:

\begin{verbatim}
try:
    data = json.loads(response.text)
    return EventThreads.model_validate(data)
except (json.JSONDecodeError, ValidationError):
    return {}  # silent fallback
\end{verbatim}

Проблема: при invalid JSON, код возвращает \texttt{{}} (пустой dict). Downstream код ожидает `EventThreads` (которая имеет required fields), получает \texttt{{}} $\to$ ValidationError $\to$ assessment молчит падает, и timeline становится пустой.

\subsubsection*{Что пробовали}
\begin{enumerate}
\item Логировать ошибку? Было, но WARNING уровень, который обрезался в production.
\item Есть ли fallback logic? Да, но fallback неполный (возвращает {} вместо default EventThreads).
\item Как это вообще прошло testing? Тесты использовали MockLLMClient, который всегда возвращает valid JSON.
\end{enumerate}

\subsubsection*{Почему не сработало}
\textbf{Exception swallowing + incomplete fallback.} Catch-all `except` absorbs both JSON errors и Validation errors. Fallback \texttt{{}} не имеет структуры, которую ожидает downstream. Это как ``promise'' вернуть результат, но фактически вернуть None.

\subsubsection*{Что сделали}
v0.7.1 (CHANGELOG, стр. 275): Resilient event\_identification added.

\begin{enumerate}
\item Distinguish между JSON parse errors и validation errors:
\begin{verbatim}
try:
    data = json.loads(response.text)
    return EventThreads.model_validate(data)
except json.JSONDecodeError as e:
    logger.error("Truncated JSON: %s", response.text[:100])
    # fallback to raw headlines?
    return _fallback_event_threads(headlines)
except ValidationError as e:
    logger.error("Validation error: %s", e)
    # return empty but structured default
    return EventThreads(threads=[])
\end{verbatim}

\item Вместо silent drop, либо:
  \begin{enumerate}
  \item Re-raise exception для explicit retry
  \item Return default объект (e.g., \texttt{EventThreads(threads=[])}) со структурой
  \item Fallback to raw headlines (e.g., каждый headline = один thread)
  \end{enumerate}
\end{enumerate}

\subsubsection*{Урок}
\textbf{Never silently swallow exceptions. Either fail-fast или graceful fallback с logged reason.}

\begin{enumerate}
\item Distinguish error types: parse errors vs validation errors имеют разные fix strategies.
\item Fallback value должен иметь корректную структуру (не \texttt{{}}, а \texttt{Model(field1=[], field2=None)}).
\item Логируй exception content (e.g. \texttt{str(e)} или \texttt{e.json()}), не только error message.
\item Unit test: add case where LLM returns truncated JSON, verify fallback works.
\end{enumerate}

---

\subsection{Case Study 13: BudgetTracker Race Condition (v0.7.1)}

\subsubsection*{Контекст}
v0.7.0 реализовала budget tracking для ограничения LLM costs (\$50 default max). BudgetTracker:

\begin{verbatim}
async def check_budget(self, cost_usd: float):
    if self.total_cost + cost_usd > self.max_budget:
        raise BudgetExceededError(...)
    self.total_cost += cost_usd
\end{verbatim}

Проблема: в async context с несколькими parallel agents, два агента могут одновременно вызвать \texttt{check\_budget()}. Race condition:
\begin{enumerate}
\item Agent 1 reads \texttt{self.total\_cost = 48}.
\item Agent 2 reads \texttt{self.total\_cost = 48}.
\item Agent 1 adds 1.5, writes \texttt{self.total\_cost = 49.5}.
\item Agent 2 adds 1.5, writes \texttt{self.total\_cost = 50.0} (мог бы быть 51, что exceeds budget).
\end{enumerate}

Результат: budget check может быть bypassed.

\subsubsection*{Что пробовали}
\begin{enumerate}
\item Запустить с parallel agents? Локально, single-threaded, не воспроизводилось.
\item Добавить sleep между LLM calls? Случайно помогло в тесте, но не fix.
\item Есть ли глобальный lock? Нет.
\end{enumerate}

\subsubsection*{Почему не сработало}
\textbf{TOCTOU (Time-of-Check-Time-of-Use) race condition.} Check и Use не атомичны. В async Pythone, `await` может переключить контекст между check и write.

\subsubsection*{Что сделали}
v0.7.1 (CHANGELOG, стр. 263):

\begin{verbatim}
class BudgetTracker:
    def __init__(self):
        self._lock = asyncio.Lock()

    async def check_budget(self, cost_usd: float):
        async with self._lock:
            if self.total_cost + cost_usd > self.max_budget:
                raise BudgetExceededError(...)
            self.total_cost += cost_usd
\end{verbatim}

Теперь check и write protected by lock. Не может быть race.

\subsubsection*{Урок}
\textbf{Shared state в async code требует synchronization primitives.}

\begin{enumerate}
\item Используй \texttt{asyncio.Lock()} для защиты read-modify-write операций.
\item Никогда не предполагай, что single-threaded code safe в async context.
\item Unit test: запусти с \texttt{asyncio.gather()} несколько concurrent calls, verify atomicity.
\item Code review: флаг на любое \texttt{self.field = ...} в async function --- может быть race condition.
\end{enumerate}

---

\subsection{Case Study 14: Timeout Cascade (v0.9.4)}

\subsubsection*{Контекст}
Дефолтный timeout агентов составлял 300 секунд. При работе с тяжёлыми промптами (Opus на сложных multi-agent медиациях) LLM-ответы регулярно превышали 300s. Результат: каскадный отказ --- таймаут одного агента вызывал падение всей стадии, а с ней терялась вся предыдущая работа pipeline.

\subsubsection*{Что сделали}
Дефолтный timeout поднят с 300s до 600s (commit c96ed67). Индивидуальные таймауты настроены per stage:

\begin{enumerate}
\item \texttt{outlet\_historian}: снижен до 300s (легковесная стадия, не требует длинных LLM-ответов).
\item \texttt{delphi\_r2}: увеличен до 900s (сложная multi-agent медиация с 5 персонами, самый длинный промпт в pipeline).
\item Остальные стадии: 600s (2x от наблюдаемого p95).
\end{enumerate}

\subsubsection*{Урок}
\textbf{Дефолтные таймауты должны базироваться на 2x от наблюдаемого p95 latency, а не на оптимистичных средних.}

\begin{enumerate}
\item Мониторь реальные latency перед выставлением порогов.
\item Настраивай таймауты per stage --- одна стадия может быть 10x тяжелее другой.
\item Каскадный отказ от таймаута хуже, чем дополнительные 5 минут ожидания.
\end{enumerate}

---

\subsection{Case Study 15: Эволюция max\_tokens (v0.5.1--v0.9.5)}

\subsubsection*{Контекст}
Параметр \texttt{max\_tokens} прошёл 4 итерации:
\begin{enumerate}
\item \textbf{4096} (начальное значение) --- обрезал выходы Delphi R1, где 5 персон генерируют по 200+ токенов каждая.
\item \textbf{8192} --- всё ещё недостаточно для полных assessment'ов с обоснованиями.
\item \textbf{16384} --- обнаружилось, что OpenRouter резервирует \texttt{max\_tokens} из кредитного баланса. Установка 16384 резервировала \$5+ за вызов, даже если фактический выход --- 2000 токенов.
\item \textbf{unlimited} (текущее значение) --- cap снят полностью.
\end{enumerate}

\subsubsection*{Что сделали}
Удалён cap \texttt{max\_tokens} полностью. Ключевое открытие --- механизм резервирования кредитов OpenRouter: при \texttt{max\_tokens=16384} система блокировала эквивалент полного output на балансе пользователя, создавая иллюзию высоких расходов и блокируя параллельные вызовы при низком балансе.

\subsubsection*{Урок}
\textbf{Модели ценообразования LLM API неочевидны.} Тестируй cost impact параметра \texttt{max\_tokens} до production deploy.

\begin{enumerate}
\item Проверяй, как провайдер обрабатывает \texttt{max\_tokens}: billing по фактическому output или по зарезервированному?
\item Для OpenRouter: unlimited \texttt{max\_tokens} дешевле, чем высокий фиксированный cap.
\item Документируй pricing gotchas для каждого LLM-провайдера в проектной wiki.
\end{enumerate}

---

\subsection{Case Study 16: Инкрементальное сохранение pipeline (v0.9.5)}

\subsubsection*{Контекст}
Pipeline сохранял результаты только по полному завершению всех 9 стадий. Если QualityGate (stage 9) падал по таймауту после 40+ минут обработки, все результаты предыдущих стадий (\$5--15 в LLM costs) терялись безвозвратно. Пользователь видел пустую страницу.

\subsubsection*{Что сделали}
Реализовано per-stage инкрементальное сохранение (commit f4d0848):

\begin{enumerate}
\item Каждая завершённая стадия немедленно persists результат в \texttt{PipelineStep.output\_data}.
\item При падении на стадии $N$, стадии $1 \ldots N{-}1$ уже сохранены и доступны пользователю.
\item Worker при рестарте может продолжить с последней сохранённой стадии (checkpoint recovery).
\end{enumerate}

\subsubsection*{Урок}
\textbf{Долгоживущие pipeline обязаны делать checkpoint после каждого дорогого шага. ``Всё или ничего'' неприемлемо, когда каждый шаг стоит доллары.}

\begin{enumerate}
\item Стоимость потерянной работы пропорциональна количеству стадий без checkpoint'ов.
\item Checkpoint recovery позволяет retry с последней точки, экономя и время, и деньги.
\item Паттерн: каждый stage записывает \texttt{output\_data} сразу после успешного завершения.
\end{enumerate}

---

\section{Отложенные направления}

\subsection{Case Study 17: Domain-Specific Brier Scores}

\subsubsection*{Идея}
Разные типы рынков имеют разные точность-характеристики. Например, crypto markets более волатильны, политические markets более точны. Можно ли compute per-domain Brier Score weights и используй их для adaptive extremizing?

\subsubsection*{Почему отложено}
\textbf{Data sparsity.} из 348K informed бетторов, только ~100 имеют >5 resolved bets в each domain (crypto, politics, sports, etc.). Brier Score на малом N имеет огромную дисперсию. При попытке stratify по domains, 99% бетторов drop out из-за недостаточного разрешения данных.

Документация (docs/methodology-inverse-problem.md, стр. 398): ``Domain-specific BS --- data sparsity (1--3\% BSS gain, not worth effort)''

\subsubsection*{Урок}
\textbf{Не добавляй complexity для marginal gains на limited data.} 1--3\% улучшение требует 10x больше данных для статистической significance.

---

\subsection{Case Study 18: Bettor-Level News Correlation}

\subsubsection*{Идея}
Корреляция между news events (GDELT) и bettor positions. Гипотеза: informed бетторы быстро реагируют на breaking news, можно ли построить predictive model?

\subsubsection*{Почему отложено}
\textbf{Спекулятивная гипотеза требует специального pipeline.} Для отладки нужно:
\begin{enumerate}
\item GDELT API, которая работает (исправлена в v0.9.4, но limited к 3-месячному окну).
\item RSS pipeline для индивидуальных outlet'ов (NewsScout есть, но не per-bettor).
\item Temporal alignment: news timestamp vs trade timestamp (сложная временная модель).
\item Validation: есть ли реальная correlation? Не known.
\end{enumerate}

Документация (docs/methodology-inverse-problem.md, стр. 399): ``Bettor-level корреляция с новостями --- требует GDELT/RSS pipeline, спекулятивная гипотеза''

\subsubsection*{Урок}
\textbf{Research hypotheses требуют pilot data перед full engineering.} Не начинай building без evidence, что ideá заслуживает effort.

---

\subsection{Case Study 19: Hierarchical Trader Belief Models}

\subsubsection*{Идея}
Bettor positions отражают их beliefs. Можно ли build hierarchical Bayesian model где бетторы cluster по beliefs, и используй это для better aggregate forecast?

\subsubsection*{Почему отложено}
\textbf{Pure research project, не engineering task.} Это требует:
\begin{enumerate}
\item Новых statistical methods (Bayesian hierarchical modeling).
\item Publication-quality validation (peer review, reproducibility).
\item Большого датасета для training (мы имеем Polymarket, но это single market type).
\item Четкого определения ``belief'' в контексте prediction markets.
\end{enumerate}

Документация (docs/methodology-inverse-problem.md, стр. 400): ``Иерархические модели --- research project, не engineering task''

\subsubsection*{Урок}
\textbf{Distinguish между engineering (solve known problem) и research (formulate new problem).} Не заполняй product roadmap research ideas.

---

\subsection{Case Study 20: Kalshi API Integration}

\subsubsection*{Идея}
Kalshi --- второй major US prediction market (после Polymarket). Может ли Delphi интегрировать Kalshi для расширения охвата?

\subsubsection*{Почему отложено}
\textbf{US-only coverage, low ROI.} Kalshi имеет:
\begin{enumerate}
\item ~500 active markets (vs Polymarket ~5000+).
\item Regional exposure (США только).
\item Higher regulation (SEC oversight).
\item No public bettor profiles (в отличие от Polymarket, который имеет on-chain transparency).
\end{enumerate}

Cost-benefit: интеграция требует ~3 дня (API integration, schema changes, tests). Результат: +200 рынков (low quality), 0 additional signals (no public bettors). Not worth.

\subsubsection*{Урок}
\textbf{Evaluate ROI перед adding new data source.} Не все APIs стоят интегрирования.

---

\subsection{Case Study 21: BigQuery for GDELT Historical Data}

\subsubsection*{Идея}
GDELT имеет полный historical dataset в BigQuery (2.65 TB/year). Можно ли использовать для batch retrospective testing?

\subsubsection*{Почему отложено}
\textbf{Cost-prohibitive at scale.} BigQuery pricing:
\begin{enumerate}
\item 1 TB free/месяц, затем \$6.25/TB.
\item Unoptimized query: 311 GB scanned = \$1.94 per query.
\item Optimized query с temporal decorators: 6.28 GB scanned = \$0.04 per query.
\item Для 100-query batch retrospective: \$400--\$200 (depends on optimization).
\end{enumerate}

Alternative: 15-minute CSV polling (free, 2--5 MB per batch) + local DuckDB. Total cost: \$0. Speed: slower, но достаточна для production monitoring.

Документация (\texttt{tasks/research/gdelt\_api.md}, стр. 43--47): ``BigQuery viable для retrospective, но для real-time --- используй CSV polling''

\subsubsection*{Урок}
\textbf{Evaluate cloud costs vs local alternatives. Free polling beats paid queries.}

---

\section*{Summary Table: Dead Ends and Lessons}

\begin{table}[h]
\centering
\small
\begin{tabular}{|l|l|l|l|l|l|}
\hline
\textbf{\#} & \textbf{Category} & \textbf{Case Study} & \textbf{Status} & \textbf{Version} & \textbf{Impact} \\
\hline
1 & \textbf{API} & Metaculus 403 & Fixed, Disabled & v0.5.1, v0.9.4 & Auth-required tier lock \\
2 &  & GDELT Cyrillic & Fixed & pre-v0.5.1, v0.9.4 & HTML crashes prevented \\
3 &  & Polymarket camelCase & Fixed & pre-v0.5.1 & 422 error resolved \\
4 &  & Reuters RSS & Removed & v0.9.4 & 404 dead feed \\
\hline
5 & \textbf{Architecture} & YandexGPT & Removed & v0.8.0 & Consolidated to OpenRouter \\
6 &  & Sonnet 4.6 & Removed & v0.9.4 & Non-existent model \\
7 &  & Dark mode & Removed & v0.8.0 & Unmaintained feature \\
8 &  & Pico.css $\to$ Tailwind & Migrated & v0.8.0 & Classless CSS limit \\
\hline
9 & \textbf{Critical Bugs} & Temporal leak & Fixed & v0.9.2 & BSS +0.092$\to$+0.127 \\
10 &  & conditionId & Fixed & v0.9.3 & 99\% signal loss \\
11 &  & Date serialization & Fixed & v0.9.5 & Complete data loss \\
12 &  & PromptParseError & Fixed & v0.5.2--v0.9.4 & Silent assessment drop \\
13 &  & BudgetTracker race & Fixed & v0.7.1 & Budget bypass \\
14 &  & Timeout cascade & Fixed & v0.9.4 & Cascading stage failures \\
15 &  & max\_tokens evolution & Fixed & v0.5.1--v0.9.5 & Credit reservation bloat \\
16 &  & Incremental save & Fixed & v0.9.5 & \$5--15 result loss \\
\hline
17 & \textbf{Deferred} & Domain-specific BS & Deferred & --- & 1--3\% gain, high sparsity \\
18 &  & Bettor-news correlation & Deferred & --- & Speculative hypothesis \\
19 &  & Hierarchical beliefs & Deferred & --- & Pure research \\
20 &  & Kalshi API & Deferred & --- & US-only, low ROI \\
21 &  & BigQuery GDELT & Deferred & --- & Cost-prohibitive (\$0.04--\$1.94/query) \\
\hline
\end{tabular}
\end{table}

\noindent
\textbf{Key Takeaways:}
\begin{enumerate}
\item \textbf{External APIs require continuous monitoring.} Breaking changes happen (Metaculus, Reuters). Subscribe to official channels.
\item \textbf{Single-provider architecture better than unreliable fallbacks.} YandexGPT was removed, OpenRouter is sole LLM provider.
\item \textbf{Serialization and temporal correctness are non-negotiable.} Date serialization crash and temporal leak both caused complete data loss / model degradation.
\item \textbf{Distinguish engineering from research.} Hierarchical models are research; bettor signals engineering.
\item \textbf{ROI analysis prevents feature creep.} Domain-specific BS, Kalshi, BigQuery: all deferred because cost > benefit.
\end{enumerate}
